\chapter{Haag--Ruelle Scattering Theory}
\label{ch:haag-ruelle}

\section{Scattering States as Hilbert-Space Limits}

The previous chapters constructed local fields, their spectral content, and
Lorentzian time-ordered Green functions. We now add a new kind of object:
asymptotic particle states. These are vectors in the physical Hilbert space
obtained as large-time limits of localized wave packets. Once these limits are
constructed, the scattering operator is a comparison between two Hilbert-space
identifications: one made in the far past and one made in the far future.

The construction in this chapter is for a massive stable bosonic particle. It
uses a Hilbert space \(\Hilb\), a vacuum vector \(\Omega\), local operators, the
Poincare representation, the spectrum condition, and an isolated
one-particle mass shell. These hypotheses are stated explicitly because the
large-time limit is a theorem under hypotheses, not an algebraic manipulation
of Green functions.

Let spacetime have dimension \(D\), and write \(d=D-1\). The mostly-plus
metric is
\[
  \eta_{\mu\nu}=\operatorname{diag}(-,+,\ldots,+),
  \qquad
  p^2=\eta_{\mu\nu}p^\mu p^\nu .
\]
Let \(U(a,\Lambda)\) be a strongly continuous unitary representation of the
proper orthochronous Poincare group on \(\Hilb\). Its translation subgroup is
written
\[
  U(a,\mathbf 1)=\ee^{\ii a_\mu P^\mu},
\]
where the commuting self-adjoint generators \(P^\mu\) are the
energy-momentum operators. The spectrum condition is
\[
  \operatorname{Spec}(P)\subset \overline V_+
  =
  \{p\in\mathbb R^D:p^0\ge 0,\ p^2\le 0\}.
\]
The invariant mass operator is
\[
  M^2=-P_\mu P^\mu .
\]

\begin{definition}[Massive one-particle sector]
Fix \(m>0\). A massive one-particle sector of mass \(m\) is a closed subspace
\(\Hilb_1\subset\Hilb\) such that:
\begin{enumerate}
  \item \(\Hilb_1\) reduces the Poincare representation;
  \item \(M^2=m^2\) on \(\Hilb_1\);
  \item the spectrum of \(P^\mu\) on \(\Hilb_1\) lies on
  \[
    \Sigma_m^+
    =
    \{(\omega_{\vec p},\vec p):\vec p\in\mathbb R^d\},
    \qquad
    \omega_{\vec p}=\sqrt{\vec p^{\,2}+m^2};
  \]
  \item there is an open interval around \(m^2\) in the spectrum of \(M^2\)
  whose only spectral point is \(m^2\).
\end{enumerate}
Let \(P_1\) denote the orthogonal projection from \(\Hilb\) onto \(\Hilb_1\).
\end{definition}

The last condition is the mass-gap form of stability used in this chapter. It
says that the one-particle shell is separated from the continuum. In a scalar
theory with a lightest stable particle, the picture is:
\[
\begin{tikzpicture}[baseline=-0.55ex,scale=0.94,>=Stealth]
  \draw[->,line width=0.65pt] (-0.2,0)--(7.2,0)
    node[right] {$M^2=-P^2$};
  \draw[line width=0.75pt] (0,0.12)--(0,-0.12);
  \node[below] at (0,-0.12) {vacuum};
  \draw[blue!75!black,line width=1.1pt] (2.15,0.32)--(2.15,-0.32);
  \node[blue!75!black,below] at (2.15,-0.32) {$m^2$};
  \draw[red!75!black,line width=2.2pt] (4.65,0)--(7.0,0);
  \node[red!75!black,below] at (4.65,-0.18) {$M_{\mathrm{thr}}^2$};
  \draw[decorate,decoration={brace,amplitude=5pt},gray!75]
    (2.35,0.55)--(4.45,0.55);
  \node[gray!75!black] at (3.5,0.95) {spectral gap};
\end{tikzpicture}
\]
Here \(M_{\mathrm{thr}}\) is the lowest invariant mass in the continuum channel
under discussion.

\section{Local Operators and One-Particle Wave Packets}

To state the large-time construction without domain distractions, use bounded
local operators. A bounded operator \(A\) is localized in a bounded spacetime
region \(O\) when it belongs to the local algebra \(\Obs(O)\). Locality means
that if \(O_1\) and \(O_2\) are spacelike separated, then
\[
  [A_1,A_2]=0,
  \qquad
  A_i\in\Obs(O_i).
\]
For unbounded point fields, this is the corresponding statement after
smearing and restricting to the common invariant domain on which the products
are defined.

The particle must be creatable from the vacuum by local operations. In the
present scalar setting this means that there is a local operator \(A\) such
that
\[
  P_1A\Omega\ne0 .
\]
Choose a Schwartz test function \(\chi\in\mathcal S(\mathbb R^D)\) whose
Fourier transform is supported in a small neighborhood of the positive mass
shell \(\Sigma_m^+\) and avoids the rest of the energy-momentum spectrum.
Define
\[
  B=A(\chi)
  =
  \int_{\mathbb R^D}\dd^D x\,\chi(x)\,A(x),
  \qquad
  A(x)=U(x)AU(x)^{-1}.
\]
The operator \(B\) is almost local: it is approximated in norm by local
operators in larger and larger bounded regions, with errors decreasing faster
than any inverse power of the size of the region. This property is sufficient
for the locality estimates used below.

Let \(h\in C_c^\infty(\mathbb R^d)\) be a smooth compactly supported momentum
wave packet. Define the positive-energy Klein--Gordon solution
\[
  h_t(\vec x)
  =
  \int\frac{\dd^d\vec p}{(2\pi)^d}\,
  h(\vec p)\,
  \ee^{\ii\omega_{\vec p}t-\ii\vec p\cdot\vec x}.
\]
The associated Haag--Ruelle operator is
\[
  B_t(h)
  =
  \int_{\mathbb R^d}\dd^d\vec x\,
  h_t(\vec x)\,B(t,\vec x),
  \qquad
  B(t,\vec x)=U(t,\vec x)BU(t,\vec x)^{-1}.
\]
The one-particle vector created by this family is
\[
  \psi_h:=P_1B_t(h)\Omega .
\]
It is independent of \(t\). The role of the phase in \(h_t\) is precisely to
cancel the free time evolution on the mass shell.

The velocity support of \(h\) is
\[
  V(h)
  =
  \left\{
    \frac{\vec p}{\omega_{\vec p}}:
    \vec p\in\operatorname{supp}h
  \right\}
  \subset\mathbb R^d .
\]
Stationary phase gives the spacetime localization of \(h_t\): for large
\(|t|\), the wave packet is concentrated near rays
\[
  \vec x=t\,\frac{\vec p}{\omega_{\vec p}},
  \qquad
  \vec p\in\operatorname{supp}h .
\]
Thus two wave packets with disjoint velocity supports become spacelike
separated as \(|t|\to\infty\).

\begin{center}
\begin{tikzpicture}[scale=0.94,>=Stealth]
  \draw[->,line width=0.7pt] (-3.35,-2.0)--(-3.35,2.45)
    node[above] {$x^0$};
  \draw[->,line width=0.7pt] (-3.35,-2.0)--(3.5,-2.0)
    node[right] {$\vec x$};
  \draw[gray!55,dashed] (-3.0,-1.55)--(2.75,1.85);
  \draw[gray!55,dashed] (-1.15,-1.55)--(3.15,1.45);
  \fill[blue!10,draw=blue!75!black,line width=1.0pt]
    (-1.85,-1.1) ellipse (0.58 and 0.34);
  \fill[green!12,draw=green!50!black,line width=1.0pt]
    (-0.6,-1.1) ellipse (0.58 and 0.34);
  \fill[blue!10,draw=blue!75!black,line width=1.0pt]
    (-0.35,0.65) ellipse (0.55 and 0.34);
  \fill[green!12,draw=green!50!black,line width=1.0pt]
    (1.75,0.35) ellipse (0.55 and 0.34);
  \draw[->,blue!75!black,very thick] (-1.65,-0.78)--(-0.62,0.35);
  \draw[->,green!50!black,very thick] (-0.35,-0.78)--(1.38,0.13);
  \node[blue!75!black] at (-1.15,1.03) {$V(h_1)$};
  \node[green!50!black] at (2.25,0.78) {$V(h_2)$};
  \draw[<->,purple!80!black,very thick] (0.22,0.62)--(1.18,0.48);
  \node[purple!80!black,align=center] at (0.73,1.18)
    {spacelike\\separation};
\end{tikzpicture}
\end{center}

\section{Multi-Particle Limits}

Let \(B_1,\ldots,B_n\) be almost-local operators obtained as above, and let
\(h_1,\ldots,h_n\in C_c^\infty(\mathbb R^d)\) have pairwise disjoint velocity
supports. Define the time-dependent vector
\[
  \Psi_t(h_1,\ldots,h_n)
  =
  B_{1,t}(h_1)\cdots B_{n,t}(h_n)\Omega .
\]

\begin{theorem}[Haag--Ruelle limits]
Under the assumptions stated above, the limits
\[
  \Psi^{\mathrm{out}}(h_1,\ldots,h_n)
  =
  \lim_{t\to+\infty}\Psi_t(h_1,\ldots,h_n),
\]
\[
  \Psi^{\mathrm{in}}(h_1,\ldots,h_n)
  =
  \lim_{t\to-\infty}\Psi_t(h_1,\ldots,h_n)
\]
exist in the norm topology of \(\Hilb\). The limits depend only on the
one-particle vectors
\[
  \psi_j=P_1B_{j,t}(h_j)\Omega\in\Hilb_1,
  \qquad j=1,\ldots,n,
\]
and not on the auxiliary operators \(B_j\) used to interpolate them.
\end{theorem}

The proof has two analytic inputs. First, disjoint velocity supports imply
that the effective localization regions of \(B_{j,t}(h_j)\) separate
spacelike at large \(|t|\). Locality then gives rapid decay of commutators
\[
  [B_{i,t}(h_i),B_{j,t}(h_j)]
\]
for \(i\ne j\). Second, spectral isolation of the mass shell controls the part
of \(B_{j,t}(h_j)\Omega\) orthogonal to \(\Hilb_1\). In Cook's method one
estimates \(\frac{\dd}{\dd t}\Psi_t\); the commutator terms and the
orthogonal spectral-error terms are integrable in \(t\), hence
\(\Psi_t\) is a Cauchy family at \(+\infty\) and at \(-\infty\).

For \(n=2\), the essential estimate is visible from the norm difference
\[
  \|\Psi_{t'}-\Psi_t\|^2
  =
  \langle\Psi_{t'}-\Psi_t,\Psi_{t'}-\Psi_t\rangle .
\]
The scalar products entering this expression are vacuum correlation functions
of four almost-local operators. At large positive or negative times, the two
velocity clusters are separated by spacelike distances growing linearly in
\(|t|\). The cluster property and locality reduce the large-time scalar
product to the product of the corresponding one-particle scalar products.

\section{Fock Structure}

Let
\[
  \mathcal F_s(\Hilb_1)
  =
  \bigoplus_{n=0}^\infty \Hilb_1^{\odot n}
\]
be the bosonic Fock space over \(\Hilb_1\), where
\(\Hilb_1^{\odot n}\) denotes the symmetric \(n\)-fold tensor product and
\(\Hilb_1^{\odot0}=\mathbb C\Omega_F\). The vector \(\Omega_F\) is the Fock
vacuum.

For wave packets with disjoint velocity supports, define
\[
  \Omega_{\mathrm{out}}
  (\psi_1\odot\cdots\odot\psi_n)
  =
  \Psi^{\mathrm{out}}(h_1,\ldots,h_n),
\]
\[
  \Omega_{\mathrm{in}}
  (\psi_1\odot\cdots\odot\psi_n)
  =
  \Psi^{\mathrm{in}}(h_1,\ldots,h_n),
\]
where \(\psi_j=P_1B_{j,t}(h_j)\Omega\). The theorem above makes these
definitions independent of the choice of interpolating operators. By
linearity and continuity they extend from such vectors to all of
\(\mathcal F_s(\Hilb_1)\).

The Fock inner product is reproduced by the Hilbert-space inner product of
asymptotic states:
\[
\begin{aligned}
&\left\langle
  \Psi^{\mathrm{out}}(\psi_1,\ldots,\psi_n),
  \Psi^{\mathrm{out}}(\varphi_1,\ldots,\varphi_m)
  \right\rangle_{\Hilb}
\\
&\qquad =
  \delta_{nm}
  \sum_{\sigma\in S_n}
  \prod_{j=1}^n
  \langle\psi_j,\varphi_{\sigma(j)}\rangle_{\Hilb_1},
\end{aligned}
\]
and the same formula holds for incoming states. Therefore
\[
  \Omega_{\mathrm{out}},
  \Omega_{\mathrm{in}}
  :
  \mathcal F_s(\Hilb_1)\longrightarrow \Hilb
\]
are isometries. Their ranges are the outgoing and incoming scattering
subspaces,
\[
  \Hilb_{\mathrm{out}}=\operatorname{Ran}\Omega_{\mathrm{out}},
  \qquad
  \Hilb_{\mathrm{in}}=\operatorname{Ran}\Omega_{\mathrm{in}} .
\]

The Poincare group acts on \(\mathcal F_s(\Hilb_1)\) by the free Fock
representation \(U_F\) induced from its action on \(\Hilb_1\). The
Haag--Ruelle wave operators intertwine the physical and free actions:
\[
  U(a,\Lambda)\Omega_{\mathrm{out/in}}
  =
  \Omega_{\mathrm{out/in}}U_F(a,\Lambda).
\]
This equation states that an asymptotic \(n\)-particle state transforms as
the corresponding symmetrized tensor product of one-particle Wigner states.

\begin{center}
\begin{tikzcd}[column sep=large,row sep=large]
  \mathcal F_s(\Hilb_1)
    \arrow[r,"\Omega_{\mathrm{in}}"]
    \arrow[d,"S"']
  &
  \Hilb_{\mathrm{in}}\subset\Hilb
  \\
  \mathcal F_s(\Hilb_1)
    \arrow[r,"\Omega_{\mathrm{out}}"']
  &
  \Hilb_{\mathrm{out}}\subset\Hilb
    \arrow[u,dashed,leftrightarrow]
\end{tikzcd}
\end{center}

\section{The Scattering Operator}

If
\[
  \Hilb_{\mathrm{in}}=\Hilb_{\mathrm{out}},
\]
then incoming scattering states evolve into outgoing scattering states in the
sector under consideration. The scattering operator on the asymptotic Fock
space is
\[
  S
  =
  \Omega_{\mathrm{out}}^*\Omega_{\mathrm{in}}
  :
  \mathcal F_s(\Hilb_1)\longrightarrow \mathcal F_s(\Hilb_1).
\]
When the incoming and outgoing ranges coincide, \(S\) is unitary. The matrix
element between an incoming \(n\)-particle wave packet
\(\psi_1\odot\cdots\odot\psi_n\) and an outgoing \(m\)-particle wave packet
\(\varphi_1\odot\cdots\odot\varphi_m\) is
\[
\begin{aligned}
&\left\langle
  \varphi_1\odot\cdots\odot\varphi_m,\,
  S(\psi_1\odot\cdots\odot\psi_n)
  \right\rangle_{\mathcal F_s(\Hilb_1)}
\\
&\qquad =
  \left\langle
  \Omega_{\mathrm{out}}
  (\varphi_1\odot\cdots\odot\varphi_m),\,
  \Omega_{\mathrm{in}}
  (\psi_1\odot\cdots\odot\psi_n)
  \right\rangle_{\Hilb}.
\end{aligned}
\]
This equation is the nonperturbative definition of the scattering matrix
elements in the massive sector. The right hand side is an inner product in the
physical Hilbert space; the left hand side records the same number in the
free asymptotic particle basis.

Asymptotic completeness is the stronger statement
\[
  \Hilb=\Hilb_{\mathrm{in}}=\Hilb_{\mathrm{out}}
\]
for the sector of finite-energy states under consideration. The definition of
\(S\) on the scattering sector requires equality of the incoming and outgoing
ranges. Full asymptotic completeness asserts that every physical state in the
sector is captured by such asymptotic particle data.

\section{Relation to Local Correlation Functions}

The construction above produces the \(S\)-operator from Hilbert-space limits.
It also identifies what a later reduction formula must compute. If
\(\psi_j=P_1B_{j,t}(h_j)\Omega\) and
\(\varphi_i=P_1C_{i,t}(g_i)\Omega\), then an \(m\)-out, \(n\)-in matrix
element is
\[
\begin{aligned}
&\left\langle
  \Omega_{\mathrm{out}}
  (\varphi_1\odot\cdots\odot\varphi_m),\,
  \Omega_{\mathrm{in}}
  (\psi_1\odot\cdots\odot\psi_n)
  \right\rangle_{\Hilb}
\\
&\quad =
  \lim_{T\to+\infty}
  \left\langle
  \Omega\left|
  C_{m,T}(g_m)^*\cdots C_{1,T}(g_1)^*
  B_{1,-T}(h_1)\cdots B_{n,-T}(h_n)
  \right|\Omega
  \right\rangle .
\end{aligned}
\]
For large \(T\), the outgoing packets are localized near future timelike
directions and the incoming packets near past timelike directions. Locality
permits these vacuum matrix elements to be reorganized into time-ordered
correlation functions up to terms that vanish in the scattering limit. The
next chapter turns this statement into LSZ reduction: one extracts the above
matrix elements from the one-particle poles of time-ordered Green functions.

\section{Scope of the Massive Construction}

The massive Haag--Ruelle construction gives a precise scattering theory for
stable particles whose one-particle mass shell is isolated and whose charge is
localizable by the operators used in the theory. It also provides the
reference point for more delicate asymptotic problems. Massless particles,
long-range gauge forces, infraparticle sectors, and theories with confinement
require additional asymptotic observables or modified state spaces. Those
phenomena are part of the same program of constructing large-time particle
data, but their hypotheses and limiting procedures differ from the massive
case treated here.

The logical output of this chapter is therefore fixed. A scattering operator
has been defined before any LSZ formula and before any perturbative
interpretation of Feynman graphs as scattering amplitudes. Perturbative
scattering calculations, when available, compute matrix elements of this
operator after the reduction theorem has connected the Hilbert-space limits to
Green functions.
